\chapter{Implementation}

\section{Python code (ETL + visualization)}
The system is implemented in Python with a clear separation of concerns:
\begin{itemize}
  \item \textbf{Sync}: periodic fetch + state tracking + change detection.
  \item \textbf{ETL}: extraction, transformation, loading into PostgreSQL.
  \item \textbf{Reporting}: Matplotlib dashboards computed from warehouse queries.
\end{itemize}

\section{Application features (API + React web application)}
This project includes a complete decision-support application on top of the warehouse:
\begin{itemize}
  \item A \textbf{FastAPI backend} that exposes analytical endpoints and runs the automatic sync + ETL loop at startup.
  \item A \textbf{web dashboard} (Vite + React) that lets users search items, plot price history, analyze equipment bonuses, and browse undervalued deals.
\end{itemize}

UI screenshots are consolidated in Chapter~7.

\subsection{Automation and scheduling}
A periodic loop runs inside the backend service. It fetches market data for the allowed servers, compares a hash of the normalized payload, and triggers ETL only when the payload changes.

\placeholderfigure{Automation Loop Timeline (every 10 minutes)}{0.92\linewidth}{5cm}

\subsection{Loading with pygramETL (concept)}
The loader uses pygramETL to manage dimensions and facts through reusable abstractions.

\begin{lstlisting}[style=python,caption={Simplified loading pattern with pygramETL (dimension ensure + fact insert)}]
# 1) Ensure dimension exists and get surrogate key
item_key = dim_item.ensure({
    "item_vnum": vnum,
    "item_name": item_name,
    "item_type": item_type,
})

# 2) Insert fact row referencing the dimension key
fact_market_transaction.insert({
    "item_key": item_key,
    "time_key": time_key,
    "transaction_type_key": trx_type_key,
    "transaction_price_yang": price_yang,
    "enhancement_level": plus,
    "transaction_timestamp": batch_timestamp,
})

dwconn.commit()
\end{lstlisting}

\subsection{Matplotlib dashboards}
Dashboards are generated by querying the warehouse (via Pandas) and plotting KPIs in Matplotlib.
\begin{itemize}
  \item Top items by average price.
  \item Distribution of undervalued deals by rating.
  \item Time-series plot of an item price history.
\end{itemize}

\placeholderfigure{Matplotlib KPI Dashboard Screenshot}{0.92\linewidth}{6cm}

\section{SQL scripts used}
\begin{itemize}
  \item Warehouse creation script: constellation schema (dimensions, facts, aggregates, indexes, views).
  \item Analytical queries used by dashboards (top items, deal ratings, price history).
\end{itemize}

\begin{table}[!ht]
\centering
\caption{Examples of SQL queries used for dashboards}
\begin{tabularx}{\linewidth}{@{}lX@{}}
\toprule
\textbf{KPI} & \textbf{Query idea (simplified)} \\
\midrule
Top expensive items & Average of \texttt{fact\_price\_history.average\_price\_yang} grouped by \texttt{dim\_item.item\_name}. \\
Deals by rating & Count of rows per \texttt{fact\_undervalued\_items.deal\_rating}. \\
Price history & Join \texttt{fact\_price\_history} with \texttt{dim\_time} and filter by item \texttt{vnum}. \\
\bottomrule
\end{tabularx}
\end{table}

